% !PRINTED 400
\weberpage{400}
%<<CONT
conjugirte Paare sind. Da man in dieser Betrachtung 1 mit 2
und 3 vertauschen kann, so folgt, dass $0, [11'], [22'], [33']$
reell und alle übrigen Doppeltangenten imaginär sind.

Nach \S.\,99 können wir leicht einen Steiner'schen Complex
bilden, der die vier reellen Doppeltangenten enthält:
\[
0[11'],\ [22'][33'],\ 1[01'],\ 1'[01],\ [23][2'3'],\ [23'][2'3],
\]
woraus noch zu sehen ist, dass dieser Complex ausser den reellen
Paaren zwei conjugirte Paare und ein conjugirtes Doppelpaar
enthält.

Dies ist also der Fall 1) des vorigen Paragraphen.

2) Es seien die sieben gegebenen Geraden:
\[
0,\ 1,\ 2,\ 3,\ 3',\ 4,\ 4',
\]
und darin $0, 1, 2$ reell, $3$ zu $3'$ und $4$ zu $4'$ conjugirt. Wir
bilden drei Complexe:
\[
\begin{aligned}
(0)&\quad 1[01],\ 2[02],\ 3[03],\ 3'[03'],\ 4[04],\ 4'[04']\\
(3)&\quad 0[03],\ 1[13],\ 2[23],\ 3'[33'],\ 4[34],\ 4'[34']\\
(3')&\quad 0[03'],\ 1[13'],\ 2[23'],\ 3[33'],\ 4[3'4],\ 4'[3'4'].
\end{aligned}
\]

Der zu (0) conjugirte Complex enthält $1, 2, 3, 3', 4, 4'$ und
ist folglich mit (0) identisch, d. h. der Complex (0) ist reell.
Daraus folgt, dass $[01], [02]$ reell, $[03]$ mit $[03']$ und $[04]$
mit $[04']$ conjugirt ist. Ferner ergiebt sich auf die gleiche
Weise, dass (3), (3') conjugirte Complexe sind, und dass also
$[33']$ reell, $[34]$ mit $[3'4']$, $[34']$ mit $[3'4]$ conjugirt imaginär
sind. Da man 0, 1, 2 beliebig vertauschen darf, und ebenso 3
mit 4, so erhalten wir folgende Zusammenstellung:

\begin{center}
\gesperrt{Reelle Doppeltangenten:}
\end{center}
\[
0,\ 1,\ 2,\ [12],\ [20],\ [01],\ [33'],\ [44'].
\]

\begin{center}
\gesperrt{Conjugirte Paare:}
\end{center}
\[
\begin{aligned}
3\ \ 4 &\quad [03][04][13][14][23][24][34][34']\\
3'\ 4' &\quad [03'][04'][13'][14'][23'][24'][3'4'][3'4].
\end{aligned}
\]

Der Steiner'sche Complex, der vier reelle Paare enthält,
ist hier
\[
0[12],\ 1[20],\ 2[01],\ [33'][44'],\ [34][3'4'],\ [34'][3'4],
\]
und er enthält noch zwei conjugirte Paare. Dies ist also der
Fall 2) des \S.\,104.
